\begin{problem}[Warmup]{(a)}
    Determine the set of right cosets of the subgroup $\langle 6 \rangle$ in $\mathbb{Z}_{20}$. Check that Lagrange's Theorem 31.2 holds in this case.
\end{problem}
\unsolved

\begin{problem}[Warmup]{(b)}
    Determine the set of right cosets of the subgroup $\langle 7 \rangle$ in $U(\mathbb{Z}_{20})$. Check that Lagrange's Theorem holds in this case.
\end{problem}
\unsolved

\begin{problem}[Warmup]{(c)}
    Determine the set of right cosets of the subgroup $\{1, -1\}$ in the group of quaternions $Q_8$. Check that Lagrange's Theorem holds in this case. (Note that we are here using the shorthand notation of Exercise 28.14 for the quaternions.)
\end{problem}
\unsolved

\begin{problem}[Warmup]{(d)}
    If $H$ is a subgroup of $G$ and two cosets of $H$ share an element, then what can you say?
\end{problem}
\unsolved

\begin{problem}[Warmup]{(e)}
    Is a coset ever a subgroup?
\end{problem}
\unsolved

\begin{problem}[Warmup]{(f)}
    If $a \in Hb$, how are the two cosets $Ha$ and $Hb$ related?
\end{problem}
\unsolved

\begin{problem}[Warmup]{(g)}
    Suppose that $G$ is an infinite group, and $H$ is a subgroup of $G$ with finitely many elements. How many distinct cosets does $H$ have?
\end{problem}
\unsolved

\begin{problem}[Warmup]{(h)}
    Suppose that $G$ is an infinite group, and $H$ is a subgroup of $G$ with finitely many cosets. How many elements does $H$ have?
\end{problem}
\unsolved

\begin{problem}{1}
    Determine explicitly the right cosets of the subgroup $\langle (124) \rangle$ in the
    group $A_4$ (see Example 31.6).
\end{problem}
\unsolved

\begin{problem}{2}
    Determine explicitly the right cosets of the subgroup $\langle (124) \rangle$ in the
    whole group $S_4$.
\end{problem}
\unsolved

\begin{problem}{3}
    Find the \emph{left} cosets of the subgroup $H = \{\iota, (12)\}$ of $S_3$. How do these
    cosets compare with the right cosets of $H$ found in Example 31.2?
\end{problem}
\unsolved

\begin{problem}{4}
    How do the left cosets of $K$ compare with the right cosets found in Example 31.3?
\end{problem}
\unsolved

\begin{problem}{5}
    If $G$ is a group of order 8 and $G$ is not cyclic, why must $x^4 = 1$ for all $x$ in
    $G$? What similar statement can you make about a group of order 27? Generalize this
    further.
\end{problem}
\unsolved

\begin{problem}{6}
    Let $G$ be a group of order $p^2$, where $p$ is prime. Show that every proper subgroup
    of $G$ is cyclic.
\end{problem}
\unsolved

\begin{problem}{7}
    Replicate the argument about the number of groups of order 6 for groups of order 10.
    You should conclude that there are essentially only two such groups; namely,
    $\mathbb{Z}_{10}$ and the dihedral group $D_5$.
\end{problem}
\unsolved

\begin{problem}{8}
    \begin{enumerate}
        \item[(a)] What are the possible orders for a subgroup of a group of order 8?
        \item[(b)] Prove that any group of order 8 must have a subgroup of order 2. (This
        is easy, and does not require that you use either Theorem 31.5 or the results of
        Exercise 9.)
    \end{enumerate}
\end{problem}
\unsolved

\begin{problem}{9}
    In this problem we analyze how many groups there are of order 8. Consider the
    following cases (with hints):

    \emph{Case 1:} There is an element of order 8.

    \emph{Case 2:} If there is no element of order 8, but there is an element $a$ of order
    4, choose an element $b$ outside $\langle a \rangle$. Now, either you can find such a
    $b$ of order 2 or else all such $b$ have order 4. If you can find such a $b$ of order
    2, then, depending on whether $a$ and $b$ commute, argue that we must obtain either
    $\mathbb{Z}_2 \times \mathbb{Z}_4$ or $D_4$. Now assume that all such $b$ have order 4.
    Argue that you must obtain the quaternions $Q_8$.

    \emph{Case 3:} Now suppose that there are no elements of order 4. So, all non-identity
    elements have order 2. First, argue that the group must be abelian. Then, pick three
    elements of order 2 and argue that you must obtain
    $\mathbb{Z}_2 \times \mathbb{Z}_2 \times \mathbb{Z}_2$.
\end{problem}
\unsolved

\begin{problem}{10}
    Show that every non-identity element of a group generates the group if and only if the
    group is cyclic of prime order.
\end{problem}
\unsolved

\begin{problem}{11}
    Let $n$ be a positive integer and $1 \le k \le n$. Consider the stabilizer subgroup
    $G_k$ of $S_n$. (See Exercise 29.6.) What does Lagrange's Theorem say about
    $[S_n : G_k]$? Demonstrate this explicitly, by placing the set of right cosets of $G_k$
    into a natural one-to-one correspondence with $\{1, 2, \ldots, n\}$.
\end{problem}
\unsolved

\begin{problem}{12}
    \begin{enumerate}
        \item[(a)] Suppose $G$ is a finite group and $a \in G$. Show that $a^{|G|} = 1$.
        \item[(b)] Now prove Fermat's Little Theorem 8.7: \emph{If $p$ is prime and
        $0 < x < p$, then $x^{p - 1} \equiv 1 \pmod{p}$.}
        \item[(c)] For any positive integer $n$, $\phi(n)$ is defined to be the number of
        positive integers less than $n$ that are also relatively prime to $n$. (This is
        called \textbf{Euler's phi function}.) So, $\phi(6) = 2$ because only 1 and 5 are
        less than 6 and also relatively prime to 6. Note that if $p$ is prime,
        $\phi(p) = p - 1$. Show that if $a$ is relatively prime to $n$, then
        \[a^{\phi(n)} \equiv 1 \pmod{n}.\]
        This generalization of Fermat's Little Theorem 8.7 is called
        \textbf{Euler's Theorem}.
    \end{enumerate}
\end{problem}
\unsolved

\begin{problem}{13}
    Refer to the definition in Exercise 12 of Euler's phi function.
    \begin{enumerate}
        \item[(a)] Compute $\phi(7)$, $\phi(20)$, and $\phi(100)$.
        \item[(b)] Use Euler's Theorem to efficiently compute
        \[2^{38} \pmod{7}, \quad 17^{18} \pmod{20}, \quad 7^{82} \pmod{100}.\]
    \end{enumerate}
\end{problem}
\unsolved

\begin{problem}{14}
    Prove that the set of (right) cosets of the subgroup $\mathbb{Z}$ in $\mathbb{R}$ can
    be placed in a one-to-one correspondence with the set of real numbers $x$, with
    $0 \le x < 1$.
\end{problem}
\unsolved

\begin{problem}{15}
    Consider the group $\mathbb{S}$, the unit circle, a multiplicative subgroup of
    $\mathbb{C}^*$. Provide a simple condition on complex numbers $\alpha, \beta$
    characterizing when the cosets $\mathbb{S}\alpha$ and $\mathbb{S}\beta$ are equal, and
    prove that it works.
\end{problem}
\unsolved
